% =============================================================================
%  Notas del Profesor — Sesión 18: Taller 2 — Resolución de problemas
%  Programación Avanzada — Universidad Panamericana
%  Compilar: pdflatex sesion18_taller_2_profesor.tex
% =============================================================================
\documentclass[12pt, oneside]{article}
\usepackage[profesor]{progavanzada}

\begin{document}

\portada{Sesión 18}{Taller 2 — Resolución de problemas}{2026}

\tableofcontents
\newpage

% =============================================================================
\section{Objetivos de la sesión}
% =============================================================================

Al finalizar la sesión el alumno será capaz de:

\begin{enumerate}
  \item Buscar el elemento más cercano en una lista hacia la izquierda
        y hacia la derecha con un recorrido lineal.
  \item Construir una cadena de salida que depende de una distancia
        calculada en tiempo de ejecución.
  \item Reconocer cuándo la fuerza bruta es la estrategia correcta
        (problema de tamaño fijo y pequeño).
  \item Usar \texttt{substr} y \texttt{stoll} para procesar dígitos
        de un número leído como cadena.
\end{enumerate}

\begin{notaprofesor}[Duración estimada]
  90 minutos.  Distribución sugerida:\\
  \phantom{00} 5 min — presentación del taller \\
  \phantom{00}10 min — \textbf{Problema 1, Parte 1:} revisión guiada en pizarrón \\
  \phantom{00}30 min — \textbf{Problema 1, Parte 2:} trabajo en equipo + envío \\
  \phantom{00} 8 min — \textbf{Problema 2, Parte 1:} revisión guiada en pizarrón \\
  \phantom{00}27 min — \textbf{Problema 2, Parte 2:} trabajo en equipo + envío \\
  \phantom{00}10 min — cierre: discusión de soluciones \\[4pt]
  Los alumnos trabajan en equipos de dos o tres personas y usan el
  cuaderno Jupyter (\texttt{sesiones/sesion18\_taller\_2.ipynb}) para
  probar localmente.
\end{notaprofesor}

% =============================================================================
\section{Dinámica del taller}
% =============================================================================

\begin{notaprofesor}[Guión general de la Parte 1]
  Para cada problema, sigue este orden en el pizarrón:
  \begin{enumerate}
    \item Lee el enunciado en voz alta; detente en cada párrafo y
          pregunta ``¿qué nos dan?'' y ``¿qué nos piden?''.
    \item Escribe la entrada del ejemplo tal como aparece; etiqueta
          cada dato (``esta línea es $n$'', ``estos son los nombres'', etc.).
    \item Traza el procesamiento paso a paso con los números concretos.
    \item Escribe la salida obtenida y verifica que coincide con el
          enunciado.
    \item Lanza la pregunta de cierre sin responderla: deja esa
          generalización para el equipo.
  \end{enumerate}
\end{notaprofesor}

\begin{notaprofesor}[Cuenta en el juez]
  Verifica que todos los alumnos tengan cuenta en el juez usado antes
  de empezar la Parte 2.
\end{notaprofesor}

\newpage

% =============================================================================
\section{Problema 1: El profesor y los nombres}
% =============================================================================

\subsection{Enunciado completo}

Dados $n$ alumnos en fila, algunos con nombre conocido y otros marcados
con \texttt{?}, responder $m$ consultas: ¿cómo llamarías al alumno en
la posición $p$?

\begin{itemize}
  \item Si la posición tiene nombre conocido: imprimir el nombre.
  \item Si el conocido más cercano está a distancia $d$ a la izquierda:
        ``\texttt{derecha} $\times d$ \texttt{de Nombre}''.
  \item Si está a la derecha: ``\texttt{izquierda} $\times d$
        \texttt{de Nombre}''.
  \item Si hay dos a igual distancia (uno a cada lado):
        ``\texttt{justo enmedio de NombreIzq y NombreDer}''.
\end{itemize}

\subsection{Parte 1 — Guía de pizarrón}

\begin{notaprofesor}[Lo que escribes en el pizarrón]
  \begin{enumerate}
    \item Dibuja la fila con posiciones numeradas del 1 al 10:
    \begin{center}
      \ttfamily
      \begin{tabular}{|c|c|c|c|c|c|c|c|c|c|}
        \hline
        A & ? & ? & B & ? & ? & ? & C & ? & ? \\
        \hline
        1 & 2 & 3 & 4 & 5 & 6 & 7 & 8 & 9 & 10 \\
        \hline
      \end{tabular}
    \end{center}
    \item \textbf{Consulta $p = 3$:} ``¿Quién es el más cercano?''
          Señala a la izquierda: A en pos.\ 1, distancia 2.
          Señala a la derecha: B en pos.\ 4, distancia 1.
          B es más cercano.  B está a la \textit{derecha} de la posición 3,
          por eso lo llamamos desde su \textit{izquierda}.
          Salida: ``\texttt{izquierda de B}''.
    \item \textbf{Consulta $p = 8$:} nombre conocido (C).
          Salida: ``\texttt{C}''.
    \item \textbf{Consulta $p = 6$:} B en pos.\ 4, distancia 2.
          C en pos.\ 8, distancia 2.  Igual → enmedio.
          Salida: ``\texttt{justo enmedio de B y C}''
          (primero el de la izquierda).
    \item \textbf{Consulta $p = 10$:} solo C a la izquierda,
          distancia 2.  No hay nadie a la derecha.
          C está a la \textit{izquierda} de la posición 10,
          entonces nos movemos hacia su \textit{derecha}.
          Salida: ``\texttt{derecha derecha de C}''.
    \item Pregunta de cierre: ``¿Cómo generalizamos esto para cualquier
          fila y cualquier consulta?'' — dejar para el equipo.
  \end{enumerate}
\end{notaprofesor}

\subsection{Análisis de la solución}

El algoritmo para cada consulta $p$:

\begin{enumerate}
  \item Si \texttt{nombres[p-1] != "?"}: imprimir el nombre y continuar.
  \item Recorrer hacia la izquierda desde $p-2$ hasta 0; el primer
        nombre conocido encontrado es el más cercano a la izquierda
        (distancia $d_L = p - 1 - i$, donde $i$ es su índice en el arreglo).
  \item Recorrer hacia la derecha desde $p$ hasta $n-1$; el primer
        nombre conocido encontrado es el más cercano a la derecha
        (distancia $d_R = i - (p-1)$).
  \item Si solo existe $d_L$: imprimir ``\texttt{derecha}'' $d_L$ veces
        + ``\texttt{de Nombre}''.
  \item Si solo existe $d_R$: imprimir ``\texttt{izquierda}'' $d_R$ veces
        + ``\texttt{de Nombre}''.
  \item Si $d_L < d_R$: imprimir ``\texttt{derecha}'' $d_L$ veces +
        ``\texttt{de NombreIzq}''.
  \item Si $d_R < d_L$: imprimir ``\texttt{izquierda}'' $d_R$ veces +
        ``\texttt{de NombreDer}''.
  \item Si $d_L = d_R$: imprimir ``\texttt{justo enmedio de NombreIzq
        y NombreDer}''.
\end{enumerate}

\subsection{Solución en C++}

\begin{verbatim}
#include <iostream>
#include <vector>
#include <string>
using namespace std;

int main() {
    int n;
    cin >> n;
    vector<string> nom(n);
    for (auto& s : nom) cin >> s;

    int m;
    cin >> m;
    while (m--) {
        int p;
        cin >> p;
        p--;  // pasar a índice 0

        if (nom[p] != "?") {
            cout << nom[p] << "\n";
            continue;
        }

        int dL = -1, dR = -1;
        string nL, nR;

        for (int i = p - 1; i >= 0; i--) {
            if (nom[i] != "?") { dL = p - i; nL = nom[i]; break; }
        }
        for (int i = p + 1; i < n; i++) {
            if (nom[i] != "?") { dR = i - p; nR = nom[i]; break; }
        }

        if (dL == -1 || (dR != -1 && dR < dL)) {
            for (int i = 0; i < dR; i++) cout << "izquierda ";
            cout << "de " << nR << "\n";
        } else if (dR == -1 || dL < dR) {
            for (int i = 0; i < dL; i++) cout << "derecha ";
            cout << "de " << nL << "\n";
        } else {
            cout << "justo enmedio de " << nL << " y " << nR << "\n";
        }
    }
}
\end{verbatim}

\begin{notaprofesor}[FAQs Problema 1]
  \textbf{¿Por qué ``izquierda de B'' y no ``derecha de A'' en la posición 3?}\\
  B está a distancia 1 y A está a distancia 2.  B es el más cercano.
  B está a la derecha de la posición 3, así que el alumno está
  a la izquierda de B.

  \medskip
  \textbf{¿Puede haber más de dos nombres a la misma distancia mínima?}\\
  No.  Como hay un solo alumno por posición, solo puede haber un nombre
  conocido a distancia $d$ a la izquierda y otro a distancia $d$ a la
  derecha.  El caso de dos equidistantes siempre es uno a cada lado.

  \medskip
  \textbf{¿Qué pasa si todos los alumnos menos uno son \texttt{?}?}\\
  El problema garantiza al menos un nombre conocido, así que siempre
  habrá al menos un $d_L$ o $d_R$ definido.  No hay caso sin respuesta.
\end{notaprofesor}

\newpage

% =============================================================================
\section{Problema 2: Once Cifras}
% =============================================================================

\subsection{Enunciado}

Dado un número de 11 cifras, partirlo en tres subcadenas contiguas y
sumar las tres partes interpretadas como enteros.  Encontrar la partición
que minimiza esa suma.

\subsection{Parte 1 — Guía de pizarrón}

\begin{notaprofesor}[Lo que escribes en el pizarrón]
  \begin{enumerate}
    \item Escribe el número en el pizarrón:
          \texttt{3\,1\,1\,2\,4\,5\,1\,0\,9\,8\,8}
          y dibuja dos ``tijeras'' moviéndose por los espacios entre
          dígitos.
    \item ``Para partir en 3, necesito 2 puntos de corte.
          Si llamo $i$ al primero y $j$ al segundo, con $1 \le i < j \le 10$,
          las tres partes son los dígitos antes de $i$, entre $i$ y $j$,
          y después de $j$.''
    \item Muestra un ejemplo concreto: $i=4$, $j=7$:
          \texttt{3112} + \texttt{451} + \texttt{0988}
          $= 3112 + 451 + 988 = 4551$.
    \item ``¿Cuántas combinaciones $(i,j)$ hay?''
          Son $\binom{10}{2} = 45$.  ``¿Podemos probarlas todas?
          Sí, 45 es un número pequeñísimo.''
    \item Muestra otra partición: $i=3$, $j=7$:
          \texttt{311} + \texttt{2451} + \texttt{0988}
          $= 311 + 2451 + 988 = 3750$.  ``Esta es menor.''
    \item Pregunta de cierre: ``¿Cómo sabes que 3750 es el mínimo
          de todas las 45?'' — dejar para el equipo.
  \end{enumerate}
\end{notaprofesor}

\begin{notaprofesor}[Nota sobre la descripción del enunciado]
  El ejemplo en el enunciado del problema muestra
  ``\texttt{3112 + 451 + 0988 = 3750}'', pero
  $3112 + 451 + 988 = 4551 \ne 3750$.
  La partición correcta que da 3750 es $i=3$, $j=7$:
  ``\texttt{311} + \texttt{2451} + \texttt{0988}''
  ($311 + 2451 + 988 = 3750$).
  El enunciado tiene un error tipográfico en la descripción del ejemplo.
  La entrada y la salida correctas (\texttt{31124510988} → \texttt{3750})
  sí son correctas.  No menciones esto hasta que los alumnos hayan
  resuelto el problema.
\end{notaprofesor}

\subsection{Análisis de la solución}

Leer el número como \texttt{string} de longitud 11.
Iterar sobre todos los pares $(i, j)$ con $1 \le i < j \le 10$:

\begin{itemize}
  \item Parte 1: \texttt{s.substr(0, i)} → $a$
  \item Parte 2: \texttt{s.substr(i, j - i)} → $b$
  \item Parte 3: \texttt{s.substr(j)} → $c$
  \item Actualizar mínimo de $a + b + c$
\end{itemize}

Los ceros iniciales (\texttt{"0988"}) se eliminan automáticamente al
convertir con \texttt{stoll}.

\begin{notaprofesor}[¿Cabe en \texttt{int}?]
  El peor caso es ``\texttt{9}'' + ``\texttt{9}'' + ``\texttt{999999999}''
  $= 9 + 9 + 999\,999\,999 = 1\,000\,000\,017$, que cabe en \texttt{int}
  (máx $\approx 2.1 \times 10^9$).  Usar \texttt{long long} es buena
  práctica defensiva.
\end{notaprofesor}

\subsection{Solución en C++}

\begin{verbatim}
#include <iostream>
#include <string>
#include <climits>
using namespace std;

int main() {
    string s;
    cin >> s;

    long long minSuma = LLONG_MAX;
    int n = (int)s.size();   // debe ser 11

    for (int i = 1; i < n - 1; i++) {
        for (int j = i + 1; j < n; j++) {
            long long a = stoll(s.substr(0, i));
            long long b = stoll(s.substr(i, j - i));
            long long c = stoll(s.substr(j));
            minSuma = min(minSuma, a + b + c);
        }
    }

    cout << minSuma << "\n";
}
\end{verbatim}

\begin{notaprofesor}[FAQs Problema 2]
  \textbf{¿Qué pasa si la cadena empieza con cero?}\\
  El enunciado dice que es un número de 11 dígitos, así que no empieza
  con cero.  Las subcadenas intermedias sí pueden empezar con cero;
  \texttt{stoll} lo maneja correctamente.

  \medskip
  \textbf{Un alumno lee el número como \texttt{long long} en vez de
  \texttt{string}.  ¿Puede extraer las partes con aritmética?}\\
  Sí, pero es mucho más complicado (divisiones y módulos con potencias
  de 10).  La solución con \texttt{string} y \texttt{substr} es
  más clara y equivalente en eficiencia para este tamaño.

  \medskip
  \textbf{¿Por qué el mínimo no siempre es la partición más equilibrada?}\\
  Intuitivamente sí: 4+4+3 dígitos distribuye bien y evita números
  grandes, por eso el mínimo tiende a estar cerca de esa distribución.
  Pero el valor de los dígitos importa; por eso hay que verificar
  todas las opciones.
\end{notaprofesor}

\newpage

% =============================================================================
\section{Materiales y recursos}
% =============================================================================

\begin{itemize}
  \item \textbf{Cuaderno:} \texttt{sesiones/sesion18\_taller\_2.ipynb}
  \item \textbf{Notas alumno:} \texttt{notas\_alumno/pdf/sesion18\_taller\_2.pdf}
\end{itemize}

\end{document}
